% Solutions for VI/04.

\chapter*{Solutions for VI/04}
\addcontentsline{toc}{chapter}{Solutions for VI/04}

\paragraph{Exercise \ref{ex:vi04-target-test}.}
For \(x\in X\), the point
\[
\varphi(x)=(f_1(x),\dots,f_m(x))
\]
lies in \(Y=V(I)\) exactly when
\[
g(f_1(x),\dots,f_m(x))=0
\]
for every \(g\in I\).  This holds for every \(x\) exactly when each
\(g(f_1,\dots,f_m)\) is the zero polynomial function on \(X\), i.e. the zero element of
\(k[X]\).

\paragraph{Exercise \ref{ex:vi04-pullback}.}
For \(f,g\in k[Y]\),
\[
\varphi^\ast(f+g)
=
(f+g)\circ\varphi
=
f\circ\varphi+g\circ\varphi,
\]
and similarly
\[
\varphi^\ast(fg)
=
(f\circ\varphi)(g\circ\varphi).
\]
Constants remain constant, so the map is a \(k\)-algebra homomorphism.

\paragraph{Exercise \ref{ex:vi04-coordinates}.}
At \(x\in X\),
\[
\varphi^\ast(\bar y_i)(x)
=
\bar y_i(\varphi(x))
=
f_i(x).
\]
Thus
\[
\varphi^\ast(\bar y_i)=f_i.
\]
Since the coordinate classes generate \(k[Y]\), their images determine the homomorphism.

\paragraph{Exercise \ref{ex:vi04-contravariance}.}
For \(h\in k[Z]\),
\[
(\psi\circ\varphi)^\ast(h)
=
h\circ\psi\circ\varphi
=
\varphi^\ast(\psi^\ast(h)).
\]
Hence
\[
(\psi\circ\varphi)^\ast
=
\varphi^\ast\circ\psi^\ast.
\]

\paragraph{Exercise \ref{ex:vi04-reconstruct}.}
Set
\[
f_i=\Phi(\bar y_i).
\]
Define
\[
\varphi_\Phi(x)=(f_1(x),\dots,f_m(x)).
\]
For every \(g\in I(Y)\),
\[
g(\bar y_1,\dots,\bar y_m)=0
\]
in \(k[Y]\).  Applying \(\Phi\) gives
\[
g(f_1,\dots,f_m)=0
\]
in \(k[X]\), so \(\varphi_\Phi(X)\subseteq Y\).  Its pullback agrees with \(\Phi\) on the
generators \(\bar y_i\), hence everywhere.

\paragraph{Exercise \ref{ex:vi04-continuity}.}
Let
\[
W=V_Y(S).
\]
Then
\[
x\in\varphi^{-1}(W)
\]
if and only if
\[
s(\varphi(x))=0
\]
for every \(s\in S\), equivalently
\[
\varphi^\ast(s)(x)=0
\]
for every \(s\in S\).  Therefore
\[
\varphi^{-1}(W)=V_X(\varphi^\ast S),
\]
which is closed.

\paragraph{Exercise \ref{ex:vi04-principal-open}.}
A point \(x\in X\) belongs to \(\varphi^{-1}(D_Y(f))\) exactly when
\[
f(\varphi(x))\neq0,
\]
equivalently
\[
(\varphi^\ast f)(x)\neq0.
\]
Thus
\[
\varphi^{-1}(D_Y(f))=D_X(\varphi^\ast f).
\]

\paragraph{Exercise \ref{ex:vi04-parabola}.}
Define
\[
\alpha(t)=(t,t^2)
\]
and
\[
\beta(x,y)=x.
\]
On the parabola \(y=x^2\),
\[
\alpha(\beta(x,y))=(x,x^2)=(x,y),
\]
while
\[
\beta(\alpha(t))=t.
\]
The pullbacks are
\[
\alpha^\ast(\bar x)=t,\qquad
\alpha^\ast(\bar y)=t^2,
\]
and
\[
\beta^\ast(t)=\bar x.
\]

\paragraph{Exercise \ref{ex:vi04-axis}.}
The pullback is
\[
i^\ast:k[x,y]\to k[t],
\qquad
x\mapsto t,\quad y\mapsto0.
\]
It is surjective because \(t\) is in the image.  Its kernel is \((y)\).  Thus
\[
k[t]\cong k[x,y]/(y),
\]
so \(i\) is the closed embedding with image \(V(y)\).

\paragraph{Exercise \ref{ex:vi04-affine-line-automorphism}.}
The inverse map is
\[
u\mapsto a^{-1}(u-b).
\]
The pullback
\[
\varphi^\ast:k[u]\to k[t]
\]
sends
\[
u\mapsto at+b.
\]
Its inverse sends
\[
t\mapsto a^{-1}(u-b).
\]

\paragraph{Exercise \ref{ex:vi04-squaring}.}
The pullback is
\[
s^\ast:k[u]\to k[t],
\qquad
u\mapsto t^2.
\]
Its image is
\[
k[t^2].
\]
If \(k\) is infinite, substitution \(u\mapsto t^2\) sends a nonzero polynomial to a
nonzero polynomial, so the map is injective.  It is not surjective because \(t\notin
k[t^2]\).

\paragraph{Exercise \ref{ex:vi04-maps-to-a1}.}
A morphism
\[
X\to\mathbb A_k^1
\]
is a single polynomial function \(f\in k[X]\).  Conversely every \(f\in k[X]\) defines
the morphism \(x\mapsto f(x)\).

\paragraph{Exercise \ref{ex:vi04-iso-ring}.}
The inverse \(k\)-algebra homomorphism
\[
(\varphi^\ast)^{-1}:k[X]\to k[Y]
\]
corresponds to a morphism
\[
\psi:Y\to X.
\]
Contravariance gives
\[
(\psi\circ\varphi)^\ast
=
\varphi^\ast\circ\psi^\ast
=
\operatorname{id}_{k[X]},
\]
and similarly for \(\varphi\circ\psi\).  By uniqueness in the affine correspondence,
\[
\psi\circ\varphi=\operatorname{id}_X,
\qquad
\varphi\circ\psi=\operatorname{id}_Y.
\]

\paragraph{Exercise \ref{ex:vi04-closed-quotient}.}
Because \(Y\subseteq X\),
\[
I(X)\subseteq I(Y).
\]
Hence restriction induces
\[
k[X]=k[x_1,\dots,x_n]/I(X)
\twoheadrightarrow
k[x_1,\dots,x_n]/I(Y)=k[Y].
\]
The kernel is
\[
I(Y)/I(X)=I_X(Y).
\]

\paragraph{Exercise \ref{ex:vi04-challenge-product}.}
If
\[
f:Z\to X,\qquad g:Z\to Y
\]
are morphisms, define
\[
(f,g):Z\to X\times Y,
\qquad
z\mapsto(f(z),g(z)).
\]
Its coordinates are simply the coordinate functions of \(f\) followed by those of \(g\),
so it is a morphism.  Conversely, composing a morphism
\[
h:Z\to X\times Y
\]
with the two coordinate projections gives a pair of morphisms to \(X\) and \(Y\).  The two
constructions are inverse.


\paragraph{Exercise \ref{ex:vi04-legacy-hyperbola}.}
The coordinate ring is
\[
k[H]=k[x,y]/(xy-1)\cong k[x,x^{-1}].
\]
This ring has nonconstant units \(cx^n\), whereas
\[
k[t]^\times=k^\times.
\]
Hence \(k[H]\not\cong k[t]\), so
\[
H\not\cong\mathbb A_k^1.
\]

A morphism
\[
\varphi:\mathbb A_k^1\to H
\]
corresponds to a homomorphism
\[
\varphi^*:k[x,x^{-1}]\to k[t].
\]
Since \(x\) is a unit, \(\varphi^*(x)\) must be a unit of \(k[t]\), hence
\[
\varphi^*(x)=c\in k^\times.
\]
Then
\[
\varphi^*(y)=c^{-1}.
\]
Thus every such morphism is constant:
\[
t\longmapsto(c,c^{-1}).
\]

\paragraph{Exercise \ref{ex:vi04-legacy-products}.}
Write
\[
X=V(I)\subseteq\mathbb A_k^m,
\qquad
Y=V(J)\subseteq\mathbb A_k^n.
\]
Inside
\[
k[x_1,\dots,x_m,y_1,\dots,y_n],
\]
extend \(I\) and \(J\) and set
\[
K=I\,k[x,y]+J\,k[x,y].
\]
A point \((x,y)\) lies in \(V(K)\) exactly when \(x\in X\) and \(y\in Y\).  Hence
\[
X\times Y=V(K).
\]
Moreover,
\[
\begin{aligned}
k[X\times Y]
&\cong k[x,y]/K\\
&\cong (k[x]/I)\otimes_k(k[y]/J)\\
&\cong k[X]\otimes_k k[Y].
\end{aligned}
\]
If \(k\) is algebraically closed and \(X,Y\) are irreducible affine varieties, their
coordinate rings are finitely generated integral \(k\)-algebras and are geometrically
integral.  Their tensor product is therefore a domain, so \(X\times Y\) is irreducible.

\paragraph{Exercise \ref{ex:vi04-legacy-cubic-map}.}
Let
\[
\varphi^*:k[x,y]/(y^2-x^3+x)\to k[t]
\]
be induced by a morphism, and put
\[
a=\varphi^*(x),\qquad b=\varphi^*(y).
\]
Then
\[
b^2=a^3-a=a(a-1)(a+1).
\]
Because \(\operatorname{char}(k)\neq2\), the three factors are pairwise coprime in the UFD
\(k[t]\).  Their product is a square, so each factor is a square up to a unit.  Since \(k\)
is algebraically closed, each nonzero unit is itself a square.  Thus for some \(p,q\in k[t]\),
\[
a=p^2,
\qquad
a-1=q^2.
\]
Subtracting,
\[
(p-q)(p+q)=1.
\]
Both factors are units of \(k[t]\), hence constants.  Therefore \(p,q\), and thus \(a\),
are constant.  Then \(b^2=a(a-1)(a+1)\) is constant, so \(b\) is constant as well.
Therefore \(\varphi\) is constant.
